% HD-R14. GPT-6 Astra Pro /root; inherited effort not exposed.
\subsection{HD-R14: actual rank-three generators and tensor decompositions}
\label{sec:rank-three-tensor-obstruction}

Let \(V\) be a complex vector space of dimension \(m\ge15\),
and let \(N\subset GL(V)\) be a finite absolutely irreducible
group generated by elements \(x\) with
\(\operatorname{rank}(x-I)=3\). There is no tensor
decomposition
\[
 \Phi:V\xrightarrow{\ \sim\ }\bigotimes_{i=1}^{k}W_i,
 \qquad k\ge2,\quad d_i=\dim W_i\ge2,                 \tag{R14.1}
\]
preserved by \(N\), including a decomposition on which
\(N\) permutes isomorphic tensor factors. The claim concerns
an actual linear isomorphism \(\Phi\) and the original
operators \(\Phi x\Phi^{-1}\). All the identifications used
in its proof are given below. In particular, a primitive
normal subgroup generated by actual rank-three matrices in
the remaining sixfold coefficient ranks is also tensor
indecomposable, as required by the primitive classification
used in HD-G31.

\paragraph{1. Exact eigenspace bound for two nonscalar factors.}
Take invertible operators \(A\) and \(B\) on spaces of
dimensions \(r,s\ge2\), with \(A\otimes B\) of finite
order. From
\(A^q\otimes B^q=I\) one obtains
\(A^q=aI\), \(B^q=a^{-1}I\) for a nonzero scalar \(a\).
For completeness, equality of two nonzero pure tensors in
\(\operatorname{End}(W)\otimes\operatorname{End}(U)\)
forces their factors to be proportional: applying a linear
functional nonzero on one factor first determines the other,
and substitution determines the first. This proves the
stated scalar identities. The polynomials \(t^q-a\) and
\(t^q-a^{-1}\) have distinct roots, so \(A,B\) are
semisimple.

Suppose both are nonscalar. Let \(a_1\ge a_2\ge\cdots\)
be the positive eigenvalue multiplicities of \(A\), and
similarly \(b_1\ge b_2\ge\cdots\) for \(B\), padding
with zeros. For every nonzero scalar \(z\), tensoring
eigenvectors gives the literal eigenspace decomposition
\[
 \ker(A\otimes B-zI)
   =\bigoplus_{\lambda\mu=z} W_\lambda\otimes U_\mu.
                                                               \tag{R14.2}
\]
The matching of distinct eigenvalues is a partial bijection.
The elementary rearrangement inequality follows by swapping
two out-of-order terms, whose change is
\((a_i-a_j)(b_i-b_j)\ge0\). Hence its dimension is at most
\(\sum_i a_i b_i\). Since both operators are nonscalar,
\(1\le a_1\le r-1\), \(1\le b_1\le s-1\), and
\[
 \begin{aligned}
 \sum_i a_i b_i
 &\le a_1b_1+(r-a_1)b_2\\
 &\le(r-1)b_1+b_2\\
 &\le(r-1)b_1+(s-b_1)\\
 &=(r-2)b_1+s
 \le(r-1)(s-1)+1.
 \end{aligned}                                               \tag{R14.3}
\]
The second inequality uses \(b_1\ge b_2\) and
\(a_1\le r-1\); the last also holds for \(r=2\).
Thus every eigenspace has codimension at least
\(r+s-2\). When \(rs=m\ge15\), one has \(r+s\ge8\):
the products of positive integer pairs with sum at most7
are at most12. Consequently an operator with both tensor
factors nonscalar has moving rank at least6.

\paragraph{2. A nontrivial tensor-factor permutation.}
Fix an operator \(x\) preserving (R14.1) and permuting its
factors. Its actual expression consists of invertible maps
\(A_i:W_i\longrightarrow W_{\pi(i)}\), followed by the
permutation of tensor slots. Isolate a cycle
\((1,\ldots,\ell)\) of length \(\ell\ge2\), whose
factors all have dimension \(d\). Choose a basis in the
first factor and transport it successively by
\(A_1,\ldots,A_{\ell-1}\). These are explicit changes of
basis. In the resulting coordinates its cycle operator is
\[
 R(u_1\otimes\cdots\otimes u_\ell)
       =P u_\ell\otimes u_1\otimes\cdots\otimes u_{\ell-1},
 \qquad P=A_\ell\cdots A_1.                              \tag{R14.4}
\]
After the literal reordering of tensor slots, the whole
operator is \(R\otimes S\), with \(S\) on the remaining
cycles (or a one-dimensional scalar factor when there are
none). Finiteness of \(x\) and the pure-tensor argument in
paragraph1 show that a power of \(R\) is scalar. Also
\(R^\ell=P\otimes\cdots\otimes P\), so a power of
\(P\) is scalar. All of \(P,R,S\) are semisimple.

Choose an eigenbasis \(e_1,\ldots,e_d\) of \(P\).
In the resulting tensor basis, \(R\) cyclically rotates
the index words and multiplies by a nonzero eigenvalue of
\(P\). On one word orbit, an eigenvector of eigenvalue
\(z\) is determined by one coefficient: successive
coefficients are fixed by the equations \(Rv=zv\).
There is therefore at most one such eigenvector per orbit.
Exactly \(d\) words are fixed by the cyclic rotation;
every other orbit has at least two words. It follows that
every eigenspace of \(R\) has dimension at most
\[
       d+\frac{d^\ell-d}{2}
       =\frac{d^\ell+d}{2}
       \le\frac34d^\ell.                                 \tag{R14.5}
\]
The last inequality is \(2d\le d^\ell\), valid for
\(d,\ell\ge2\). Applying (R14.2) to the semisimple
operators \(R,S\) bounds every eigenspace of \(x\) by
\(3m/4\). In particular
\[
                   \operatorname{rank}(x-I)\ge m/4>3.
                                                               \tag{R14.6}
\]
Thus every rank-three generator induces the identity
permutation of the tensor factors. Because those matrices
generate \(N\), the entire group fixes every factor in
the decomposition.

\paragraph{3. The preserved-factor case gives a proper invariant subspace.}
Write a rank-three generator in the coordinates (R14.1) as
\(\bigotimes_i A_i\). If two factors are nonscalar,
separate one from the tensor product of all the others.
The latter is nonscalar by the pure-tensor equality
argument. Paragraph1 then gives moving rank at least6,
contradicting its actual rank3. If every factor is scalar,
the operator is scalar on \(V\), and its moving rank is
zero or \(m\), again impossible. There is exactly one
nonscalar factor, indexed by \(i\).

Absorb the product of all other scalar factors into that
one actual operator, denoting the result by \(A_i'\).
This leaves the original matrix unchanged and gives
\[
 x=A_i'\otimes I_{\widehat W_i},\qquad
 \widehat W_i=\bigotimes_{j\ne i}W_j,\qquad
 3=\operatorname{rank}(A_i'-I)\dim\widehat W_i.           \tag{R14.7}
\]
Since \(k\ge2\), the cofactor dimension is at least2.
Thus it is exactly3, the active moving rank is1, and the
active factor has dimension \(m/3\ge5\). A product of
integers at least2 equal to3 has just one factor; hence
there are exactly two tensor factors, of dimensions
\(m/3\) and3. The factor of dimension3 cannot itself be
active for a rank-three generator, because its cofactor
dimension is at least5. Every generator therefore acts
identically on the same three-dimensional factor, in the
unchanged expression (R14.7).

Choose a nonzero vector \(u\) in that factor. The actual
linear map
\[
 W_i\longrightarrow V,\qquad
                  w\longmapsto\Phi^{-1}(w\otimes u)        \tag{R14.8}
\]
is injective, has proper nonzero image, and intertwines
each generator with its active operator \(A_i'\).
The image is consequently invariant under \(N\),
contradicting absolute irreducibility. This proves the
assertion, including tensor permutations.

For an original coefficient representation over an
imaginary quadratic \(K\), apply the argument to its
literal complexification \(V_K\otimes_K\mathbb C\).
It proves a nonexistence statement about a tensor
isomorphism of that representation. No new coefficient
field, abstract replacement character, or unproved
polarized transport is introduced; all original curve
and Jacobian maps remain those of HD-S44 and HD-G31.
